% Pursuing Stacks (cotes 134-2 to 134-8): archive cote 134-5, site page 104 (PDF page 105), typescript page 291.
% typescript. Draft compared: ps5.tex:1-64.
% AI-transcribed by Claude Opus 5.5 (claude-opus-5-5, adaptive thinking, effort high), Message Batches API, batch ps2 (msgbatch_01NVwr1Yuo8qJf5rykeGBXPK), 2026-11-12; not yet checked by eye.
\scanpage{134-5/104}
\mspage{291}

1.7.

87.    I could'nt resist last night and had to look through Thomason's
preprint on the closed model structure of (Cat).  The paper is  really
pleasant reading  -  and it gives exactly what had been la\unsure{kk|ck}ing me in my
reflections lately on the homotopy theory of (Cat)  -  namely a class of
neat monomorphisms  $Y\hand{\longrightarrow}X$  which \inserted{all} have the property that cobase change
by\xout{\unsure{x}} these preserves weak equivalence,  the so-called \uline{Dwyer} \uline{maps}.  I had
hoped for a while that``open immersions''  and their duals,  the ``closed
immersions''  (namely sieve and cosieve maps,  in Thomason's wording), have
this property,  and when it turned out they had'nt,  I had been at a loss
of what stronger property put in instead,  wide enough however to allow
for the standard factorisation statements for a map to go through.  The
definition of a Dwyer map is an extremely pretty one  -  it is an open
immersion $Y\hand{\to}X$  such that the induced map from Y into it's closure \hand{$\overline{Y}$}
should have a right adjoint.  Now this implies that $Y\hand{\to}\overline{Y}$  is aspheric,
and I suspect that this \inserted{extra} condition on an open immersion $Y\hand{\longrightarrow}X$ should be
sufficient,  to imply that it has the ``cofibration property'' above with
respect to weak equivalence.  It would mean in a sense that the given
open immersion is very close to be a closed immersion too,  without how-%
ever being a direct summand necessarily.  The dual notion  is that of a
\margin{\hand{$\mathring{Y}$} (circled)}closed immersion such that the corresponding interior  $\hand{\mathring{\phantom{Y}}}\!\!Y$  of Y in X gives
a map  $\mathring{Y}\hand{\longrightarrow}Y$  which has a left adjoint,  or \inserted{only} which is ``coaspheric'' .
By \xout{\unsure{the}} Qui\over{ll}{ll}en's duality principle,  if one notion works well for push-%
outs,  so does the other.  -  With this notion in hands,  it should'nt
be difficult now to get a closed model structure on (Cat) a lot simpler
than Thomason's.  Visibly he was hampered by the standing reflex :
homotopy = semi-simplicial algebra,  which caused him to pass by the de-%
tour of the category $\hand{\Delta}\hat{}$ of semi-simplicial complexes,  rather than just
working in (Cat) itself.  I'll \inserted{have to} come back upon this in part V  of the notes\hand{,}
where I intend to investigate the homotopy properties of (Cat) and elemen-%
tary modelizers $A\hat{}$,  including the existence of closed model structures.

    Some comments of Thomason's at the end of his preprint,  about
applications to algebraic K-theory,  seem to indicate that the notion
of ``integration'' and ``cointegration'' of homotopy types I have been inte-%
rested in,  has been studied (under the name holim and \over{\unsure{c}}{h}ocolim) in the
context of closed model categories by Anderson  (his paper appeared in
1978).   As I am going to develop some ideas along these lines in part VI
on derivators,  I should have a look at what Anderson does,  notably

\margin{104}
